\documentclass[12pt,a4paper]{article}

\usepackage[margin=1in]{geometry}
\usepackage{hyperref}
\usepackage{graphicx}
\usepackage{booktabs}
\usepackage{enumitem}
\usepackage{microtype}
\usepackage{setspace}
\usepackage{fancyhdr}
\usepackage{titlesec}
\usepackage{amsmath}
\usepackage{float}

\onehalfspacing

\title{
    \vspace{2cm}
    \textbf{\Huge MultiModal Mobility Platform} \\
    \vspace{0.5cm}
    \Large Unifying Fragmented Territories through Data Integration \\
    \vspace{2cm}
}
\author{
    \textbf{Master 1 IoT} \\
    Muhammad Omais Rana \\
    \vspace{0.5cm}
    \textbf{Supervisor:} Philippe Canalda \\
    \vspace{1cm}
}
\date{\today}

\begin{document}

\maketitle
\thispagestyle{empty}

\newpage
\tableofcontents
\newpage

\section{Introduction \& Context}
The subject of this project is to conceptualize, design, and implement a unified, multi-modal Mobility as a Service (MaaS) platform. This platform is specifically tailored for regions where mobility services are heavily fragmented across various disconnected operators. By bridging these gaps, the platform aims to provide end-users with seamless, door-to-door journey planning that dynamically integrates buses, trains, and shared vehicles, ultimately encouraging a massive shift towards greener transportation.

\section{Problematics \& Challenges}

\subsection{Fragmented Territory}
One of the most significant barriers to seamless public transportation is the geographic and administrative fragmentation of the territory. In regions like Bourgogne-Franche-Comté, the mobility landscape resembles a "mille-feuille" (layer cake) of Autorités Organisatrices de la Mobilité (AOMs) and Opérateurs de Mobilité (OMs). Each administrative block, whether a municipality, an agglomeration, or a broader regional council, manages its own specific transport network. 

Because these networks are designed in silos, their service areas often end abruptly at administrative borders, creating severe discontinuities for passengers attempting to travel between neighboring towns.

\begin{figure}[H]
    \centering
    \includegraphics[width=0.8\textwidth]{mille.png}
    \caption{Map illustrating the fragmented boundaries of various AOMs and OMs.}
    \label{fig:fragmented_map}
\end{figure}

To resolve this, the platform calculates overlapping transfer points, essentially stitching these disjointed regions together to form a continuous travel corridor.

\begin{figure}[H]
    \centering
    % \includegraphics[width=0.8\textwidth]{mille_bridged.png}
    \vspace{4cm} % Placeholder space
    \caption{Diagram illustrating how fragmented territories and disparate transport networks are successfully bridged via multimodal hubs.}
    \label{fig:mille_bridged}
\end{figure}

\subsection{Aggregation of Heterogeneous Data}
Solving the fragmentation problem requires the aggregation of massive, \textbf{heterogeneous, multimodal data types}. A core tenet of this project is that the platform operates on \textbf{real data, not simulations}. Every route generated relies on actual production data published by transit authorities.

This data is fundamentally split into distinct paradigms:
\begin{itemize}[leftmargin=2em]
    \item \textbf{Static Data (GTFS):} Traditional public transit operates on predetermined schedules. This data is universally encoded using the General Transit Feed Specification (GTFS). GTFS is highly structured but inherently static, relying on massive sets of CSV files defining precise geographic stops, polyline shapes, and calendar-specific service exceptions.
    \item \textbf{Dynamic Data (GBFS):} Modern micromobility (shared bikes, scooters, and cars) operates on a free-floating basis. This data is represented using the General Bikeshare Feed Specification (GBFS), which utilizes lightweight JSON APIs to broadcast live vehicle locations, battery levels, and availability states.
    \item \textbf{Carpooling Data:} A third emerging data type involves peer-to-peer ad-hoc journeys, requiring a dynamic registry of driver commutes that must be reconciled against the fixed GTFS schedules.
\end{itemize}

The architectural challenge lies in building a pipeline capable of parsing gigabytes of static GTFS schedules in SQL databases while continuously polling live GBFS endpoints, and ultimately merging these datasets into a unified routing graph\footnote{Future implementations will also evaluate real-time transit standards such as GTFS-RT, SIRI, and NeTEx to provide live delay notifications and vehicle positions.}.



\newpage

\section{Proposed Solution Architecture}

\subsection{Global Infrastructure \& Data Flow}
The solution is built on a robust multi-tenant architecture. The underlying infrastructure separates the hardware layers (database clusters, background queue workers) from the software application layer. The data flow acts as a Directed Acyclic Graph (DAG) where raw GTFS archives and live GBFS streams are ingested, normalized into a common relational schema, and then indexed for rapid spatial queries.

\begin{figure}[h!]
    \centering
    % \includegraphics[width=0.9\textwidth]{infrastructure_dag.png}
    \vspace{4cm} % Placeholder space
    \caption{Global Diagram detailing hardware/software infrastructure and the directional flow of data from external sources to the end-user.}
    \label{fig:infrastructure_dag}
\end{figure}

\subsection{Actor Authorities \& Multi-Tenancy}
The architecture isolates data by cryptographically tagging every transport element with a \texttt{tenant\_id}. The system identifies three major actors, each wielding distinct authorities:
\begin{enumerate}
    \item \textbf{AOM (Superadmin):} The regional authority that possesses a global overview of the entire territory. The AOM is primarily responsible for the \textbf{management of territories and all subordinate OMs}. Crucially, the AOM also dictates high-level \textbf{Pricing power} across the network and analyzes aggregated \textbf{customer usage statistics} to evaluate the performance of the entire region.
    \item \textbf{OM (Tenant Admin):} The specific mobility operator. They maintain full autonomy over their specific service lines, possessing the authority to upload their own GTFS schedules, manage employee roles, and track tenant-specific performance without relying on the AOM.
    \item \textbf{End User (Customer/Guest):} Interacts exclusively with the unified search portal to consume the aggregated data.
\end{enumerate}

\begin{figure}[h!]
    \centering
    % \includegraphics[width=0.9\textwidth]{actors_authorities.png}
    \vspace{4cm} % Placeholder space
    \caption{Interaction diagram mapping the three major actors to their respective authorities and UI panels.}
    \label{fig:actors_authorities}
\end{figure}

\begin{figure}[H]
    \centering
    % \includegraphics[width=0.9\textwidth]{aom_dashboard.png}
    \vspace{4cm} % Placeholder space
    \caption{Screenshot of the AOM (Regional Administrator) dashboard interface, displaying overarching regional statistics and multi-tenant management features.}
    \label{fig:aom_dashboard}
\end{figure}

\begin{figure}[H]
    \centering
    % \includegraphics[width=0.9\textwidth]{om_dashboard.png}
    \vspace{4cm} % Placeholder space
    \caption{Screenshot of the OM (Mobility Operator) tenant dashboard, highlighting the GTFS upload module and localized fleet analytics.}
    \label{fig:om_dashboard}
\end{figure}

\begin{figure}[H]
    \centering
    % \includegraphics[width=0.9\textwidth]{end_user_interface.png}
    \vspace{4cm} % Placeholder space
    \caption{Screenshot of the End User public interface, showing the unified multimodal map and custom itinerary search capabilities.}
    \label{fig:end_user_interface}
\end{figure}

\newpage

\section{Routing Algorithms \& Data Processing}

\subsection{Implementation Tech Stack}
The backend is driven by \textbf{Laravel 12 (PHP 8.3)}, utilizing its expressive ORM and robust queuing system for background data ingestion. The administration panels are constructed using \textbf{Filament v4}, providing a rapid Server-Driven UI. A highly relational \textbf{MySQL} database relies on spatial indexing to quickly calculate Haversine distances between coordinates and transport nodes.

\subsection{Custom Heuristic Routing Algorithm}
Industry-standard transit algorithms like A* or Dijkstra rely on graph-theory. Converting a time-dependent GTFS network into a graph results in a massive "state explosion" that is highly inefficient to traverse within a standard relational database. Conversely, timetable-based algorithms like RAPTOR (Round-Based Public Transit Routing) avoid graphs but require the entire regional dataset to be pre-loaded into RAM—a significant architectural challenge for stateless PHP applications handling multiple tenants.

Instead of fighting these architectural limitations, this platform leverages a deterministic, data-driven heuristic optimized for relational SQL databases.

\begin{enumerate}[label=\textbf{\arabic*.}, leftmargin=*]
    \item \textbf{Spatial Resolution:} The algorithm locates the nearest bus stops or shared vehicles within a defined walking radius from the user's origin.
    \item \textbf{Trunk and Feeder Leg Construction:} 
    \begin{itemize}
        \item It identifies major hubs (e.g., SNCF stations) near both the origin and destination.
        \item A "Feeder" bus is located to transport the user to the initial hub.
        \item A "Trunk" train is located to traverse the massive distance between the fragmented territories.
        \item A final "Feeder" vehicle completes the journey to the workplace.
    \end{itemize}
    \item \textbf{CO2 Emission Calculation:} As the algorithm compiles the valid legs of the journey, it applies a carbon coefficient based on the transport mode (e.g., $X$ grams/km for an electric train vs $Y$ grams/km for a diesel bus), providing the user with an exact ecological footprint for the chosen itinerary.
    \item \textbf{GBFS Injection:} Live shared bicycles are evaluated as direct alternatives for short distances, scored against public transit based on total duration.
\end{enumerate}

\begin{figure}[h!]
    \centering
    % \includegraphics[width=0.8\textwidth]{itinerary_flow.png}
    \vspace{4cm} % Placeholder space
    \caption{Flow diagram detailing the step-by-step logic of the itinerary generation algorithm.}
    \label{fig:itinerary_flow}
\end{figure}

\begin{figure}[h!]
    \centering
    % \includegraphics[width=0.8\textwidth]{algo_screenshot.png}
    \vspace{4cm} % Placeholder space
    \caption{Code snippet demonstrating the execution of the custom routing algorithm.}
    \label{fig:algo_screenshot}
\end{figure}

\newpage

\section{Deliverables \& Results}

To ensure the long-term viability and academic reproducibility of this project, the following assets and contributions have been produced:

\begin{itemize}[leftmargin=2em]
    \item \textbf{Platform (GitHub):} The entire source code of the backend, database migrations, and frontend map interfaces are version-controlled and available for further development.
    \item \textbf{Demonstration Environment:} A fully functional live instance allowing users to perform multimodal searches crossing the boundaries of multiple real-world transit networks.
    \item \textbf{Technical Documentation:} Comprehensive User's and Developer's guides covering MaaS navigation, GTFS ingestion, and local Docker configurations.
    \item \textbf{Scientific Research Paper (Under Preparation):} An academic paper synthesizing the findings, methodological frameworks, and real-world impact of this multi-tenant routing architecture. This publication is intended to serve as the foundational research for a potential Ph.D. thesis exploring advanced Mobility as a Service paradigms.
    \item \textbf{Conference Contribution:} Results and methodological frameworks from this project are being prepared for presentation at an upcoming mobility/IoT conference meeting.
\end{itemize}

\section{Future Perspectives}

\subsection{Carpooling Integration (MuCar)}
While public transit forms the backbone of the platform, private vehicles remain ubiquitous in rural areas devoid of transit stops. A primary future objective is the integration of the "MuCar" module—a peer-to-peer carpooling network. Private drivers will register their commutes directly into the routing engine, allowing the algorithm to dynamically suggest a citizen's car as a "Feeder" leg to transport a passenger to the nearest major train station.

\begin{figure}[h!]
    \centering
    % \includegraphics[width=0.8\textwidth]{mucar_integration.png}
    \vspace{4cm} % Placeholder space
    \caption{Conceptual diagram illustrating the integration of the MuCar carpooling network into the existing GTFS/GBFS routing engine.}
    \label{fig:mucar_diagram}
\end{figure}

\subsection{Unified Pricing and Ticketing}
Achieving a true "Level 3" MaaS platform requires transitioning from route discovery to financial settlement. Future iterations will introduce global pricing bundles orchestrated by the AOM. A user will pay a single monthly subscription to the platform, and a complex financial clearinghouse algorithm will track validation metrics to fairly distribute the revenue back to the specific OMs.

\newpage

\section{Conclusion}
The geographic and administrative fragmentation of mobility services poses a critical barrier to sustainable urban and rural transportation. This project successfully demonstrates that a centralized, software-driven architecture can overcome these physical borders without stripping autonomy from individual mobility operators. By adopting a strict multi-tenant infrastructure and utilizing real, heterogeneous data streams—combining static GTFS schedules with dynamic GBFS APIs—the platform actively weaves disparate networks into a unified travel corridor. 

Through the implementation of a custom, CO2-aware routing algorithm that prioritizes privacy, end-users are provided with seamless multimodal itineraries that actively encourage a shift away from individual thermal vehicles. Ultimately, the integration of public transit, private micromobility, and future peer-to-peer carpooling within a single ecosystem lays a definitive technical foundation for the future of Mobility as a Service in fragmented regions.

\newpage

\section{References}

\begin{thebibliography}{9}

\bibitem{datagouv}
Data.gouv.fr, \textit{Plateforme ouverte des données publiques françaises}. Transport and Mobility datasets. Available at: \url{https://transport.data.gouv.fr/}

\bibitem{gtfs}
MobilityData, \textit{General Transit Feed Specification (GTFS) Reference}. Detailed documentation on the static transit modeling standard.

\bibitem{gbfs}
MobilityData, \textit{General Bikeshare Feed Specification (GBFS)}. Documentation for dynamic micromobility status representation.

\bibitem{maas_models}
Hensher, D. A. (2017). \textit{Future bus transport contracts under a mobility as a service (MaaS) regime in the digital age: Are they likely to change?}. Transportation Research Part A: Policy and Practice. Provides extensive reading on mobility models and public-private transitions.

\bibitem{routing_algo}
Bast, H., et al. (2016). \textit{Route Planning in Transportation Networks}. Defines the standard complexities of algorithm design for time-expanded and time-dependent transit networks.

\end{thebibliography}

\end{document}
